\section{Discussion}
\label{sec:discussion}

In this section we synthesize the experimental findings to discuss the implications of BRS for operating system design and its practical value in real-world settings \cite{Bouron2018Battle,10306992,6365626}.

\subsection{Rethinking the Responsiveness--Fairness Trade-off}
Our results show that the widely held assumption of an inherent fairness--responsiveness trade-off does not always hold \cite{10306992,4663063,8360522}: through carefully targeted changes to the scheduler's decision logic, \textbf{BRS} improves responsiveness with only bounded fairness deviation. This reframes a decades-old question---whether to prioritize throughput or perceived latency \cite{Bouron2018Battle,241774,8149743}---by replacing a strict fairness principle with a model that can be tuned to diverse QoS requirements \cite{5580094,5169938,8338088}, addressing the severe resource contention of modern mobile systems that conventional static schedulers handle poorly \cite{9121657}.

\subsection{Implications for OS Design and Research}
The broader lesson is that \emph{small, well-chosen kernel changes can reshape system behavior} \cite{Bouron2018Battle,6365626,4663063}: anchoring responsiveness control in the kernel frees applications from working around scheduling blind spots \cite{10306992,917709,6919315}; a \textit{vruntime} bias lets one general-purpose scheduler balance latency, fairness, and energy jointly \cite{8360522,9979771,1244943}; and kernel-level placement makes the OS workload-aware rather than treating tasks as interchangeable units \cite{6919315,1336761}.

\subsection{Industrial and Production Considerations}
BRS is most valuable where latency is a safety or financial risk: it lets real-time cloud and edge systems meet SLA targets without a specialized real-time kernel, and its predictable timing benefits AI-robotics control loops and AR/VR jitter. The conservative defaults ($\alpha{=}0.2,\beta{=}0.15$) suit most environments; aggressive tuning should be reserved for specific latency-sensitive systems \cite{10306992,7426385,5580094}.


\subsection{Limitations and Open Challenges}
While BRS offers several strengths, several areas remain open. Beyond the x86\_64 evaluation (Section~\ref{sec:evaluation}), the precise interactivity-score and bias-bound definitions (Section~\ref{subsec:vruntime-nudge}), the strengthened stability analysis (Section~\ref{subsec:math-stability}), and the adversarial-robustness study (Section~\ref{subsec:adversarial}) added in this version~\cite{8758850,8894365}, highly specialized workloads such as high-frequency transactions, avionics, and constrained MCU-based systems warrant dedicated study~\cite{9999639,10007654}.

\textbf{Heterogeneous and CPU--GPU platforms.} BRS currently spans ARM big.LITTLE and x86\_64, where schedulable entities are symmetric CPU threads and core asymmetry is a performance ratio. On CPU--GPU integrated SoCs or AMD-style big+small designs, two issues arise that bounded \textit{vruntime} bias does not by itself resolve. First, performance asymmetry means a fixed bias $\alpha B_i$ yields different \emph{absolute} latency gains on a big versus a LITTLE core, so the coefficient should be made core-class aware (e.g.\ scaled by capacity), which interacts non-trivially with the stability analysis. Second, GPU work is dispatched through a separate command-queue model that CPU \textit{vruntime} accounting does not observe, so end-to-end responsiveness for CPU--GPU pipelines requires coordinating BRS with the GPU driver's scheduler. We expect the bound of Definition~\ref{def:bound} still holds per CPU core, but uniform tail-latency gains across asymmetric cores will require capacity-aware logic rather than parameter tuning alone~\cite{6882237,6899614,9979771}---the most important next step for generality.

Finally, integrating machine-learning-based prediction beyond reactive bias---to enable proactive rather than corrective adjustment---is a promising direction for deeper energy-efficiency gains in dynamic environments~\cite{7426385,Zhang2025Tail}.


\subsection{Operational Tuning Scenarios}
For deployment we distill three recurring scenarios, each handled by a small \texttt{/proc} adjustment that the hybrid controller then maintains automatically. \emph{(A)~Latency spikes during frame bursts}---briefly raise the tie-breaker with a short TTL (\texttt{echo 30 > \$BRS/bias\_beta\_boost\_ttl\_seconds}), auto-reverted. \emph{(B)~Fairness-floor breach ($J<0.96$)}---dampen the bias (\texttt{echo 0.02 > \$BRS/bias\_alpha\_decrement}) and let hybrid mode recover. \emph{(C)~Background starvation $>2\%$}---the auto-clamp engages on its own; verify with \texttt{cat \$BRS/stats}. Here \texttt{\$BRS} is \texttt{/proc/sys/kernel/sched\_brs}. These small, easily overlooked adjustments are the practical safeguards that keep the system stable as load shifts~\cite{5580094,5169938}.

\subsection{Future Directions}
\label{subsec:future}
We plan to extend BRS in three directions \cite{10306992,6919315}. \textbf{Adaptive and predictive control.} We are evolving BRS into a self-tuning system in which a lean controller monitors live signals---runqueue spikes or interrupt-request (IRQ) bursts---and nudges $(\alpha,\beta)$ to stay within a safe, SLO-meeting range rather than chasing peak performance \cite{7426385,Zhang2025Tail,6782279,8338088}. \textbf{Cross-layer scalability.} We are integrating BRS with energy-efficient scheduling and the Dynamic Voltage and Frequency Scaling (DVFS) governor on heterogeneous hardware, and adding hierarchical \texttt{cgroup} and Non-Uniform Memory Access (NUMA)-aware placement so it scales from large x86 clusters to low-power ARM cores \cite{1244943,9121657,6919315,8360522}. \textbf{Formal verification.} We are modeling the responsiveness--fairness frontier with closed-loop stability and using coverage-guided fuzzing (\texttt{syzkaller}) to scrutinize safety properties under fault injection \cite{5169938,9351534}.
